% ─────────────────────────────────────────────────────────────────────────────
\section{Results}
\label{sec:results}
% ─────────────────────────────────────────────────────────────────────────────

All results in this section are taken directly from the project notebooks
(\texttt{03\_AD\_Models}, \texttt{04\_FC\_Models}, \texttt{05\_E2E\_Models}) and
the metric tables they write to \texttt{results/}.

\subsection{Anomaly Detection}
\label{subsec:results_ad}

Table~\ref{tab:ad_results} reports 5-fold cross-validation performance for the
classical models on Repr.~B (184 features), the tuned Ensemble on
Repr.~F\textsubscript{++} (224 features), and the published DL
benchmarks.\footnote{All deep learning figures in this report are reproduced
from the LMSD preprint \cite{lmsd2025}. That preprint was subsequently
\emph{withdrawn} by its authors, who stated that ``significant methodological
flaws have been identified in the experimental validation and metric-computation
procedures.'' We therefore present these numbers as \emph{indicative} reference
points from the literature, not as validated upper bounds, and we interpret the
classical-vs.-DL gaps accordingly throughout.}

\begin{table}[H]
    \centering
    \caption{Anomaly Detection results (5-fold CV). Best classical result in
             \textbf{bold}; best DL result in \textit{italics}.}
    \label{tab:ad_results}
    \begin{tabular}{llcccc}
        \toprule
        \textbf{Model} & \textbf{Repr.} & \textbf{F1} & \textbf{AUC} & \textbf{FNR} & \textbf{FPR} \\
        \midrule
        \multicolumn{6}{l}{\textit{Classical models (this work, Repr.~B)}} \\
        Ensemble       & B (184) & 0.691 & 0.754 & 0.322 & 0.296 \\
        Logistic Reg.  & B (184) & 0.679 & 0.730 & 0.308 & 0.333 \\
        XGBoost        & B (184) & 0.675 & 0.742 & 0.330 & 0.321 \\
        Linear SVM     & B (184) & 0.675 & 0.726 & 0.347 & 0.303 \\
        Random Forest  & B (184) & 0.651 & 0.707 & 0.400 & 0.298 \\
        KNN            & B (184) & 0.545 & 0.562 & 0.477 & 0.431 \\
        \midrule
        \textbf{Ensemble (best)} & \textbf{F\textsubscript{++} (224)} & \textbf{0.703} & \textbf{0.767} & \textbf{0.307} & \textbf{0.287} \\
        \midrule
        \multicolumn{6}{l}{\textit{Deep learning benchmarks (indicative) \cite{lmsd2025}}} \\
        \textit{ConvTokMHSA}  & raw & \textit{0.799} & --- & \textit{0.163} & --- \\
        InceptionTime \cite{inceptiontime2020} & raw & 0.784 & --- & 0.196 & --- \\
        MMK Net        & raw & 0.721 & --- & 0.298 & --- \\
        \bottomrule
    \end{tabular}
\end{table}

\subsubsection{Analysis}

\paragraph{Best classical model.}
The tuned soft-voting Ensemble on Repr.~F\textsubscript{++} achieves the best
classical result, F1\,=\,0.703 with AUC\,=\,0.767 and FNR\,=\,0.307. On the
plain Repr.~B feature space the same Ensemble reaches F1\,=\,0.691, narrowly
ahead of Logistic Regression (0.679), XGBoost (0.675), and Linear SVM (0.675).
KNN trails badly (0.545), confirming that distance-based classification degrades
in the 184-dimensional feature space. The near-balanced AD task ($\approx$1.04:1)
means F1-macro tracks binary F1 closely, and a linear or shallow model on global
statistics is already a competitive health screen.

\paragraph{Cross-physical gain.}
Moving from Repr.~B to Repr.~F\textsubscript{++} (adding inter-cylinder CHT/EGT
difference features) lifts the Ensemble from F1\,=\,0.691 to 0.703 ($+$1.2\,pp)
and lowers FNR from 0.322 to 0.307. This confirms that inter-cylinder temperature
asymmetry carries information beyond per-sensor global statistics.

\paragraph{Gap to DL.}
The best DL model (ConvTokMHSA, F1\,=\,0.799) outperforms the best classical
system by \textbf{9.6\,pp in F1} and \textbf{14.4\,pp in FNR}; our Ensemble is,
however, within \textbf{1.8\,pp} of MMK~Net (0.721). ConvTokMHSA applies global
self-attention over the full 2{,}048-step sequence, allowing it to detect subtle
whole-flight coherence patterns---such as gradual thermal drift correlated with
electrical anomalies across flight phases---that are invisible to any global
statistic. This gap is the indicative cost of discarding temporal structure for
AD.

\paragraph{Feature importance.}
The tuned Ensemble's feature importances are dominated by oil-pressure statistics:
the top channels are \texttt{E1 OilP\_std}, \texttt{E1 OilP\_range}, and
\texttt{E1 OilP\_mean}, followed by \texttt{E1 CHT1\_max}, \texttt{E1 RPM\_std},
and electrical features (\texttt{amp2\_std}). Oil pressure being the single most
discriminative channel is consistent with its role as the primary engine-health
indicator and direct sensor for lubrication-system faults.

\subsection{Fault Classification}
\label{subsec:results_fc}

Table~\ref{tab:fc_results} reports FC results on the 5{,}602 anomalous flights,
compared to the DL benchmarks.

\begin{table}[H]
    \centering
    \caption{Fault Classification results (5-fold CV, F1-macro). Best classical
             in \textbf{bold}.}
    \label{tab:fc_results}
    \begin{tabular}{llcc}
        \toprule
        \textbf{Model} & \textbf{Repr.} & \textbf{F1 macro} & \textbf{Accuracy} \\
        \midrule
        \multicolumn{4}{l}{\textit{Classical (this work)}} \\
        XGBoost (default)        & B (184)   & 0.116$\pm$0.007 & 0.372 \\
        XGBoost (tuned)          & B (184)   & 0.176$\pm$0.016 & 0.224 \\
        \textbf{XGBoost (tuned)} & \textbf{F\textsubscript{++} (224)} & \textbf{0.207$\pm$0.024} & \textbf{0.255} \\
        XGB cascade K\,=\,4      & F\textsubscript{++} (224) & 0.172$\pm$0.023 & --- \\
        \textbf{Oracle K\,=\,4 (ceiling)} & \textbf{F\textsubscript{++} (224)} & \textbf{0.459$\pm$0.027} & --- \\
        \midrule
        \multicolumn{4}{l}{\textit{Deep learning benchmarks (indicative) \cite{lmsd2025}}} \\
        \textit{MMK Net} & raw & \textit{0.623} & --- \\
        ConvTokMHSA      & raw & 0.208 & --- \\
        \bottomrule
    \end{tabular}
\end{table}

\subsubsection{Analysis}

\paragraph{Long-tailed difficulty.}
FC is dramatically harder than AD. The default XGBoost on Repr.~B reaches only
F1\,=\,0.116, predicting with reasonable accuracy just the two head classes
(intake gasket and rocker cover, together $\approx$55\% of anomalous flights)
while collapsing the minority classes. The per-class breakdown (notebook~04)
shows only the two head classes clearing F1\,$>$\,0.30, most classes landing in a
$0.10$--$0.30$ partial band, and per-class F1 correlating directly with
training-set size---the signature of a 38:1 long-tailed distribution rather than
a modeling defect.

\paragraph{Tuning and cross-physical features.}
Hyperparameter tuning lifts XGBoost on Repr.~B from F1\,=\,0.116 to 0.176, and
moving to Repr.~F\textsubscript{++} (cross-cylinder features) raises the best
flat classical result to F1\,=\,0.207. Inter-cylinder asymmetry is precisely the
signal needed to distinguish \emph{which} cylinder/baffle component is at fault,
so it helps FC more than the near-saturated AD task.

\paragraph{Oracle hierarchical ceiling.}
Grouping the 19 classes into 4 physically coherent super-groups via Ward
clustering and training a per-group XGBoost raises the ceiling to F1\,=\,0.459
when the coarse group is known at test time (Oracle K\,=\,4). The realistic
two-stage cascade, which must \emph{predict} the group, reaches only F1\,=\,0.172
---below the flat 0.207---showing that the hierarchical gain is entirely
contingent on correct group routing and is lost when the router is imperfect.

\paragraph{Gap to DL.}
The best DL model (MMK~Net, F1\,=\,0.623) exceeds even the Oracle ceiling
(0.459) by \textbf{16\,pp}, and the flat classical best (0.207) by 41\,pp.
Tellingly, the global-attention ConvTokMHSA (0.208) barely matches our flat
XGBoost on Repr.~F\textsubscript{++}, confirming that whole-flight context does
\emph{not} help fault typing: the discriminative signal is local, and only
locally restricted temporal models (micro-kernel convolutions) capture it.

\subsection{End-to-End Diagnosis}
\label{subsec:results_e2e}

Composing both stages into a single 20-class decision
(0\,=\,healthy, 1--19\,=\,fault type), Table~\ref{tab:e2e_results} compares the
two-stage cascade against a direct 20-class classifier and the DL reference.

\begin{table}[H]
    \centering
    \caption{End-to-end diagnosis (F1-macro over 20 classes).}
    \label{tab:e2e_results}
    \begin{tabular}{llcc}
        \toprule
        \textbf{Architecture} & \textbf{Repr.} & \textbf{F1$_{20}$} & \textbf{AD recall} \\
        \midrule
        Direct 20-class            & F\textsubscript{++} & 0.113$\pm$0.012 & 0.509 \\
        \textbf{Cascade AD$\to$FC} & \textbf{F\textsubscript{++}} & \textbf{0.131$\pm$0.026} & 0.688 \\
        \midrule
        \textit{LMSD (ConvTok$\to$MMK) \cite{lmsd2025}} & raw & \textit{0.623} & 0.837 \\
        \bottomrule
    \end{tabular}
\end{table}

The cascade AD$\to$FC (F1\,=\,0.131, AD recall 0.688) outperforms the direct
20-class classifier (F1\,=\,0.113, AD recall 0.509): decomposing the problem
keeps the strong AD stage from being diluted by the dominant \texttt{healthy}
class. Both classical pipelines remain far below the indicative LMSD figure
(0.623), but---as in FC---that gap is inherited from the fault-typing stage and
is measured against a withdrawn benchmark.

\subsubsection{Error Decomposition}

To locate \emph{where} the cascade fails, every prediction is sorted into one of
four mutually exclusive categories (Table~\ref{tab:error_decomp}).

\begin{table}[H]
    \centering
    \caption{E2E error decomposition for the cascade AD$\to$FC (notebook~05).}
    \label{tab:error_decomp}
    \begin{tabular}{lr}
        \toprule
        \textbf{Category} & \textbf{\% of predictions} \\
        \midrule
        Correct                                      & 50.0\% \\
        Stage~1 false negative (fault~$\to$~healthy) & 15.3\% \\
        Stage~1 false positive (healthy~$\to$~fault) & 14.7\% \\
        Stage~2 wrong fault type                     & 20.0\% \\
        \midrule
        \textbf{Total errors} & \textbf{50.0\%} \\
        \bottomrule
    \end{tabular}
\end{table}

Two findings follow directly:

\begin{enumerate}
    \item \textbf{Fault typing is the dominant error source.} Stage~2 wrong-type
          errors ($20.0\%$) are the single largest error category---larger than
          either Stage~1 failure mode. The anomaly is correctly detected but the
          wrong fault type is assigned, which misdirects maintenance without
          endangering the aircraft. Improving FC, not AD, is therefore the
          highest-leverage intervention.

    \item \textbf{Stage~1 false negatives are the safety-critical mode.} The
          $15.3\%$ of predictions where a faulty flight is labeled healthy never
          reach the fault classifier and are the most dangerous operational
          error; reducing FNR remains a safety priority even though it is not the
          largest error bucket.
\end{enumerate}

\subsection{Summary Comparison}
\label{subsec:summary_compare}

Table~\ref{tab:final_summary} consolidates the best result per task against the
indicative DL benchmarks.

\begin{table}[H]
    \centering
    \caption{Final performance summary. $\uparrow$\,=\,higher is better;
             $\downarrow$\,=\,lower is better. DL figures are indicative.}
    \label{tab:final_summary}
    \begin{tabular}{llccc}
        \toprule
        \textbf{Task} & \textbf{Best model / method} & \textbf{F1}$\uparrow$ & \textbf{FNR}$\downarrow$ & \textbf{vs.\ DL best} \\
        \midrule
        AD  & Ensemble Repr.~F\textsubscript{++}  & 0.703 & 0.307 & $-$0.096 \\
        AD  & ConvTokMHSA (DL)        & 0.799 & 0.163 & --- \\
        \midrule
        FC  & XGB Repr.~F\textsubscript{++} (flat) & 0.207 & --- & $-$0.416 \\
        FC  & Oracle K\,=\,4 (ceiling) & 0.459 & ---  & $-$0.164 \\
        FC  & MMK Net (DL)            & 0.623 & ---   & --- \\
        \midrule
        E2E & Cascade AD$\to$FC Repr.~F\textsubscript{++} & 0.131 & --- & $-$0.492 \\
        E2E & LMSD (DL)              & 0.623 & ---   & --- \\
        \bottomrule
    \end{tabular}
\end{table}

The results support the hypothesis stated in Section~\ref{sec:intro}: classical
models with global statistical features are \textit{competitive but not
equivalent} to DL for AD ($-$9.6\,pp vs.\ the best, within $1.8$\,pp of MMK~Net),
and face a substantially larger gap for FC ($-$16\,pp even at the Oracle ceiling,
measured against indicative DL figures). The most promising path to closing the
remaining FC gap appears to be temporal deep learning models---specifically,
locally restricted architectures like MMK~Net---rather than further feature
engineering on global statistics. We stress that the magnitude of this gap rests
on benchmark numbers from a withdrawn preprint and should be re-validated once
corrected results become available.
